Existence. Let $A$ be associative unital and $f : V \to A$ linear with $f(v)^{2} = Q(v)\,1$. By the universal property of the tensor algebra, $f$ extends to a unique unital homomorphism $T f : T(V) \to A$. It vanishes on each generator $v \otimes v - Q(v)\,1$ of $I_Q$, hence on $I_Q$, so it factors through a unique homomorphism $\bar f : T(V)/I_Q \to A$ with $\bar f \circ \iota = f$. In the quotient $\iota(v)^{2} = Q(v)\,1$ by construction.

Uniqueness. If $(C, \iota)$ and $(C', \iota')$ are both universal, the universal properties give $\phi : C \to C'$ and $\psi : C' \to C$ with $\phi \iota = \iota'$ and $\psi \iota' = \iota$. Then $\psi \phi$ and the identity are both homomorphisms $C \to C$ extending $\iota$, so they are equal by uniqueness in the universal property; likewise $\phi \psi = 1$.

The representation. With $B$ the polar form and $\iota_{\alpha}$ the derivation of $\Lambda V$ extending a linear form $\alpha$, set $c(v) = v \wedge {} + \iota_{B(v, \cdot)}$. Since $v \wedge v = 0$, $\iota_{\alpha}^{2} = 0$ and $\iota_{\alpha}(v \wedge \omega) + v \wedge \iota_{\alpha}\omega = \alpha(v)\,\omega$, we get $c(v)^{2} = B(v, v) = Q(v)$, by \claimref{std-algebra.polarisation}, so the universal property gives a homomorphism $\Cl(V, Q) \to \operatorname{End} \Lambda V$. Nothing here needs $G$ rational; the check uses rational $G$ because it computes exactly.
